% =====================================================================================
% Appendix A · The data-generating models.
%
% The equations below are the SAME equations as the chapters' §N.2 sections — they are
% definitions of the simulated worlds, not results, and are therefore the one class of
% literal this book permits in its source (see "A note on the numbers"). Every *planted
% truth* is a macro, injected from the executed notebook, exactly as in the chapters.
%
% Chapters 3 (uplift) and 8 (media mix) are cross-referenced through the \ifcsname guard
% below, so that this appendix typesets whether or not those chapters have landed: the
% guard falls back to the plain chapter number, which is what \cref would print anyway.
% =====================================================================================
\chapter{The data-generating models}
\label{app:dgp}

\newcommand{\refuplift}{\ifcsname r@ch:uplift\endcsname\cref{ch:uplift}\else Chapter~3\fi}
\newcommand{\refmmm}{\ifcsname r@ch:mmm\endcsname\cref{ch:mmm}\else Chapter~8\fi}

A causal method cannot be graded on real data, because on real data the counterfactual is missing
--- that is the entire problem. So every method in this book is first run against a world whose
answer is known by construction, and only then trusted on a world whose answer is not. The claim
each chapter then makes --- \emph{this estimator recovered the truth to within so much; that one
did not; this interval covered on so many panels out of a hundred} --- is only as auditable as the
truth it was graded against.

This appendix is that audit. It collects, in one place and as display mathematics, every world the
book plants. Nothing here is new: each model is the chapter's own \S N.2, and the chapter is where
the modelling choices are argued for. What this appendix adds is the ability to read all thirteen
worlds at once, and to check, without taking anything on trust, exactly what each estimator was
marked against.

Two conventions hold throughout. Every simulator is \textbf{seeded}, and the seed is fixed in the
notebook, so a reader who runs the code gets the panel this book analysed and not a cousin of it.
And the constants below are \textbf{definitions}, not measurements: they are the only literals this
book's source permits. Every \emph{planted truth} --- the quantity the estimators are graded
against --- is a macro, injected from the notebook that computed it, like every other number in
these pages.

% =====================================================================================
\section{Potential outcomes, confounding, and DAGs}
\label{app:dgp:fnd}

\Cref{sec:fnd:dgm} builds three worlds, because the chapter has three jobs: to show confounding,
to ground an estimator ladder, and to manufacture a collider.

\textbf{The customer file.} Each of \nbZeroNCustomers{} customers carries features $x$ ---
recency $R$, frequency $F$, monetary value $M$, tenure, engagement $E$ --- a baseline spend
$\mu_0(x)$ and a planted heterogeneous effect $\tau(x)$:
%
\begin{equation}
  Y \;=\; \mu_0(x) \;+\; \tau(x)\,T \;+\; \varepsilon,
  \qquad \varepsilon \sim \mathcal N\!\left(0,\ 8^{2}\right).
  \label{eq:app:fnd_uplift}
\end{equation}
%
The simulator is run twice and \emph{only the assignment rule changes}:
%
\begin{align}
  \text{randomized:}\quad    & T \sim \text{Bernoulli}(\tfrac12), \label{eq:app:fnd_rct}\\
  \text{observational:}\quad & T \sim \text{Bernoulli}\Big(\sigma\big(1.8\,(E - 0.4)
     + 1.2\,\tfrac{150 - R}{150} + 0.15\,(F - 5)\big)\Big). \label{eq:app:fnd_obs}
\end{align}
%
The planted average effect is \eur{\nbZeroTrueAte} in \emph{both} worlds. Nothing about the effect
changes between \cref{eq:app:fnd_rct} and \cref{eq:app:fnd_obs}; only the mechanism that assigns it
does, and that alone is enough to break the naive estimator.

\textbf{The one-confounder world}, on which the estimator ladder of \cref{sec:fnd:classical} and
the posterior of \cref{ch:pvp} are both graded. Loyalty $L$ is observed and confounds; the trait
$G$ is unobserved, moves spend, and never touches assignment --- so it is \emph{not} a confounder:
%
\begin{align}
  L,\ G &\;\sim\; \mathcal N(0, 1), \qquad
  T \;\sim\; \text{Bernoulli}\big(\sigma(0.8\,L)\big), \label{eq:app:fnd_dag_t}\\
  Y &\;=\; 20 \;+\; 5\,L \;+\; 6\,T \;+\; 3\,G \;+\; \varepsilon,
       \qquad \varepsilon \sim \mathcal N\!\left(0,\ 4^{2}\right). \label{eq:app:fnd_dag_y}
\end{align}
%
The planted effect is \eur{\nbZeroTrueDag}. The coefficient $5$ on $L$ in
\cref{eq:app:fnd_dag_y} and the coefficient $0.8$ on $L$ in \cref{eq:app:fnd_dag_t} are the two
halves of the confounding; remove either and the naive estimator would be fine.

\textbf{The zero-effect world}, in which any effect an estimator reports is manufactured rather
than merely mismeasured:
%
\begin{equation}
  T \sim \text{Bernoulli}(\tfrac12), \quad
  Y = 20 + 0 \cdot T + \varepsilon, \quad \varepsilon \sim \mathcal N(0, 5^2), \quad
  C = \mathbf 1\big\{\,1.2\,T + 1.2\,\tilde Y + \eta > 0.8\,\big\},
  \label{eq:app:fnd_collider}
\end{equation}
%
with $\eta \sim \mathcal N(0,1)$ and $\tilde Y$ spend standardized. $C$ --- ``responded to the
campaign'' --- is a common effect of the treatment and the outcome. \Cref{eq:app:fnd_collider} is
the smallest machine that manufactures a causal effect out of nothing.

\begin{remark}
\Cref{ch:pvp} plants no world of its own. It fits \crefrange{eq:app:fnd_dag_t}{eq:app:fnd_dag_y}
twice --- once by ordinary least squares, once as a posterior --- which is precisely what makes its
comparison a comparison.
\end{remark}

% =====================================================================================
\section{Uplift targeting}
\label{app:dgp:uplift}

The world of \refuplift{} is a customer file with a known, deliberately awkward,
\emph{heterogeneous} effect --- because the chapter's question is not \emph{does the discount
work?} but \emph{whom should it be sent to?}, and those have different answers.

Customers carry recency $R \sim U(5, 330)$, frequency $F \sim \text{Poisson}(5)$, monetary value
$M \sim U(15, 180)$, tenure $\sim U(1,60)$ and engagement $E \sim \text{Beta}(2,3)$. Baseline spend
is linear in them; the \emph{planted per-customer effect} is not, and its shape is the whole point
of the chapter:
%
\begin{align}
  \mu_0(x) &\;=\; \big[\,20 + 1.4\,F + 0.25\,M + 12\,E - 0.03\,R\,\big]_{+}, \label{eq:app:up_mu}\\
  \tau(x)  &\;=\; \underbrace{42\, e^{-\left(\frac{R - 120}{70}\right)^{2}}
                   \,\sigma(F - 3)\,\sigma\big(9(E - 0.28)\big)}_{\text{the persuadables}}
              \;-\; \underbrace{24\,\sigma\big(9(0.22 - E)\big)\,\sigma\big(1.2(F - 6)\big)}
                   _{\text{the sleeping dogs}}, \label{eq:app:up_tau}\\
  Y &\;=\; \mu_0(x) \;+\; \tau(x)\,T \;+\; \varepsilon,
       \qquad \varepsilon \sim \mathcal N\!\left(0,\ 8^{2}\right). \label{eq:app:up_y}
\end{align}
%
\Cref{eq:app:up_tau} is planted to be \emph{non-monotone and negative in places}. A band of
mid-recency, engaged, frequent buyers responds strongly; a pocket of disengaged frequent buyers is
\emph{harmed} by the discount --- the sleeping dogs, whom an average-effect analysis cannot see
and a targeting policy must not email. Assignment is either a coin flip or the same targeted
propensity as \cref{eq:app:fnd_obs}, and an optional hidden confounder $U$ enters both the
propensity logit and spend, with a single knob scaling each channel --- which is what the
sensitivity analysis sweeps.

% =====================================================================================
\section{Segment effects}
\label{app:dgp:seg}

\Cref{sec:seg:dgm} draws \nbTwoN{} customers with a value tier $v_i \in \{0,1\}$, an engagement
score $g_i \sim U(0,1)$, a background trend covariate $t_i$, and an email $e_i$ assigned by a fair
coin. The planted per-customer effect is linear in the two moderators and in nothing else:
%
\begin{align}
  \tau_i &\;=\; 3 \;+\; 9\,v_i \;+\; 8\,g_i, \label{eq:app:seg_tau}\\
  Y_i &\;=\; \beta_0 \;+\; \tau_i\,e_i \;+\; \beta_2 v_i \;+\; \beta_5 g_i \;+\; \beta_6 t_i
             \;+\; \varepsilon_i,
       \qquad \varepsilon_i \sim \mathcal N\!\left(0,\ 6^{2}\right). \label{eq:app:seg_y}
\end{align}
%
At the average engagement $g = 0.5$ the two tier effects are \eur{\nbTwoTrueLow} (low) and
\eur{\nbTwoTrueHigh} (high) --- \emph{those} are the grades --- and the population average effect
is \eur{\nbTwoTrueAte}. The contact cost, \eur{\nbTwoCost}, is planted \emph{between} the two, so
that the pooling choice changes what the firm does rather than merely how it reports.

% =====================================================================================
\section{Price elasticity}
\label{app:dgp:pe}

\Cref{sec:pe:dgm} produces a region-week panel: $R = \nbThreeNRegions$ regions over
\nbThreeNWeeks{} weeks, of which three are deliberately thinned so that shrinkage has something to
do (\nbThreeNObs{} region-weeks in all). Each region's elasticity is drawn once and then fixed:
%
\begin{align}
  \beta_r &\sim \mathcal N\!\left(-1.4,\ 0.35^{2}\right), \label{eq:app:pe_beta}\\
  s_t &= 0.3 \sin\!\Big(\tfrac{2\pi t}{26}\Big), \quad g_t = 0.01\,t, \quad
        C_{rt} \sim \mathcal N\!\left(20,\ 2^{2}\right), \quad
        u_{rt} \sim \mathcal N(0,1), \label{eq:app:pe_latent}\\
  \log P_{rt} &= \log 15 \;-\; 0.05\,\frac{C_{rt}}{20} \;+\; \kappa\, u_{rt} \;+\; \nu_{rt},
        \label{eq:app:pe_price}\\
  \log Q_{rt} &= 2.5 \;+\; \beta_r \log P_{rt} \;+\; 0.4\,\frac{C_{rt} - 20}{20} \;+\; s_t \;+\; g_t
        \;+\; u_{rt} \;+\; \varepsilon_{rt}, \label{eq:app:pe_demand}
\end{align}
%
with $\nu_{rt} \sim \mathcal N(0,\ 0.05^{2})$ and
$\varepsilon_{rt} \sim \mathcal N(0,\ 0.08^{2})$.

The competitor's price $C_{rt}$ appears in \emph{both} \cref{eq:app:pe_price} and
\cref{eq:app:pe_demand}, which is exactly what makes it a confounder; season and trend appear only
in demand, which makes them precision covariates rather than confounders. $\kappa$ is the
\textbf{endogeneity switch}: at $\kappa = 0$ the same demand shock $u_{rt}$ is zeroed out and
recovery is clean; turned on, it enters price and demand at once, and no pooling scheme can undo
the simultaneity. The fleet-mean elasticity is $\nbThreeTrueFleet$ and the between-region spread is
$\nbThreeTrueTau$.

% =====================================================================================
\section{Funnel mediation}
\label{app:dgp:med}

\Cref{sec:med:dgm} draws \nbFourNObs{} users through a linear structural chain, one equation per
arrow of the funnel. The onboarding score $O$ and the acquisition channel quality $Q$ are
exogenous:
%
\begin{align}
  O,\ Q &\;\sim\; U(0,1), \label{eq:app:med_exog}\\
  \text{engagement} &\;=\; 2.2\,O \;+\; 1.1\,Q \;+\; \varepsilon_1,
        &\varepsilon_1 &\sim \mathcal N(0,\ 0.4^{2}), \label{eq:app:med_eng}\\
  \text{activated} &\;=\; 1.6\,\text{engagement} \;+\; 0.6\,Q \;+\; \varepsilon_2,
        &\varepsilon_2 &\sim \mathcal N(0,\ 0.4^{2}), \label{eq:app:med_act}\\
  \text{converted} &\;=\; 0.9\,\text{activated} \;+\; 0.35\,\text{engagement} \;+\; \varepsilon_3,
        &\varepsilon_3 &\sim \mathcal N(0,\ 0.3^{2}). \label{eq:app:med_conv}
\end{align}
%
Two plants carry the chapter. There is \textbf{no $O \to \text{converted}$ arrow}, so the true
natural direct effect is exactly $\nbFourTrueNde$ and every unit of onboarding's value travels
through the funnel; because the system is linear, each route's effect is the product of its arrows'
coefficients, giving $\nbFourTrueViaAct$ through activation and $\nbFourTrueViaEng$ through
engagement alone, for a total of $\nbFourTrueTotal$. And \textbf{$Q$ feeds two stages}
(\cref{eq:app:med_eng} and \cref{eq:app:med_act}), so it is a common cause of two mediators: fail
to adjust for it and every path product built on the engagement-to-activation edge is poisoned.

% =====================================================================================
\section{What to control for}
\label{app:dgp:ctrl}

\Cref{sec:ctrl:dgm} draws \nbFiveN{} customers. The world is a superset of
\crefrange{eq:app:fnd_dag_t}{eq:app:fnd_dag_y} --- the same confounder, the same planted
\eur{\nbFiveTrueAte} --- with two \emph{post-treatment} columns added, because the chapter is about
the variables an analyst is tempted to put in the model and must not:
%
\begin{align}
  T &\;\sim\; \text{Bernoulli}\big(\sigma(0.8\,L)\big),
      \qquad L,\ G \sim \mathcal N(0,1), \label{eq:app:ctrl_t}\\
  Y &\;=\; 20 \;+\; 5\,L \;+\; 6\,T \;+\; 3\,G \;+\; \varepsilon,
      \qquad \varepsilon \sim \mathcal N(0,\,4^{2}), \label{eq:app:ctrl_y}\\
  \text{opened\_email} &\;\sim\; \text{Bernoulli}\big(\sigma(1.5\,T + 0.7\,G)\big),
      \label{eq:app:ctrl_open}\\
  \text{responded} &\;\sim\; \text{Bernoulli}\big(\sigma(1.2\,T + 0.05\,(Y - \bar Y))\big).
      \label{eq:app:ctrl_resp}
\end{align}
%
\Cref{eq:app:ctrl_resp} is a collider on $T$ and $Y$ --- caused by both, causing neither. And
\cref{eq:app:ctrl_open} is the subtler trap: a collider on $T$ and the \emph{unobserved} $G$, so a
graph drawn from the observed columns alone cannot see that it is one. Two auxiliary worlds of
\nbFiveMiniN{} customers each isolate a mediator and the M-bias structure.

% =====================================================================================
\section{Incrementality and the media mix}
\label{app:dgp:mmm}

The world of \refmmm{} is the only one in the book whose simulator has the \emph{same functional
form as the model that will be fitted to it} --- adstock carry-over and a Michaelis--Menten saturation
curve --- which is what makes its failures instructive rather than trivial. Over 156 weeks, with
$t$ the week index and $\text{season}_t = \sin(2\pi t / 52)$:
%
\begin{align}
  \text{tv}_t &\;=\; \big[\,40 + 8\,\text{season}_t + \zeta_t\,\big]_{+},
      \qquad \zeta_t \sim \mathcal N(0, 8^{2}), \label{eq:app:mmm_tv}\\
  \text{search}_t &\;=\; \big[\,25 + 11\,\text{season}_t + \xi_t\,\big]_{+},
      \qquad \xi_t \sim \mathcal N(0, 7^{2}), \label{eq:app:mmm_search}\\
  A_c(t) &\;=\; x_{c,t} \;+\; \lambda_c\, A_c(t-1),
      \qquad f_c(a) \;=\; \frac{\alpha_c\, a}{\kappa_c + a},
      \label{eq:app:mmm_adstock}\\
  \text{sales}_t &\;=\; 100 \;+\; f_{\text{tv}}\big(A_{\text{tv}}(t)\big)
       \;+\; f_{\text{search}}\big(A_{\text{search}}(t)\big)
       \;+\; \underbrace{35\,\text{season}_t}_{\text{the confounder}}
       \;+\; \varepsilon_t, \label{eq:app:mmm_sales}
\end{align}
%
with $\varepsilon_t \sim \mathcal N(0, 4^{2})$, carry-over rates $\lambda_{\text{tv}} = 0.4$ and
$\lambda_{\text{search}} = 0.2$, and saturation constants
$(\alpha, \kappa)_{\text{tv}} = (60, 50)$ against
$(\alpha, \kappa)_{\text{search}} = (8, 40)$.
%
Seasonality is the confounder and it is planted to be one twice over: it lifts sales directly
(the $35\,\text{season}_t$ term of \cref{eq:app:mmm_sales}) \emph{and} it drives spend --- brand
search far more strongly than television (\cref{eq:app:mmm_search} against
\cref{eq:app:mmm_tv}). So a naive attribution that omits it over-credits the channel that merely
\emph{follows} demand, while adjusting for it --- which is what the backdoor criterion prescribes
--- corrects the ranking. The saturation curves of \cref{eq:app:mmm_adstock} are what make the
channels' true returns computable: television has a large real effect and brand search a small one,
and both are known exactly, week by week, which is what lets the chapter grade a media-mix model's
\emph{levels} rather than merely admire its fit.

% =====================================================================================
\section{Geo lift}
\label{app:dgp:geo}

\Cref{sec:sc:dgm} produces a weekly sales panel of \nbSevenNDmas{} markets over
\nbSevenNWeeks{} weeks. Three latent factors are \emph{shared by every market} --- a secular trend,
a seasonal cycle, and a macroeconomic random walk:
%
\begin{align}
  g_t &= 0.4\,t, \qquad
  s_t = 8\sin\!\Big(\tfrac{2\pi t}{26}\Big) + 4\sin\!\Big(\tfrac{2\pi t}{13}\Big), \label{eq:app:geo_gs}\\
  m_t &= \sum_{u \le t} \eta_u, \qquad \eta_u \sim \mathcal N\!\left(0,\ 1.2^{2}\right),
        \label{eq:app:geo_macro}\\
  Y_{it} &= c_i \;+\; \lambda_{i1}\,g_t \;+\; \lambda_{i2}\,s_t \;+\; \lambda_{i3}\,m_t
           \;+\; \varepsilon_{it},
  \qquad \varepsilon_{it} \sim \mathcal N\!\left(0,\ 3^{2}\right), \label{eq:app:geo_y}
\end{align}
%
with market baselines $c_i \sim U(80, 140)$ and factor loadings $\lambda_{ij} \sim U(0.6, 1.4)$.
The campaign is planted in market $0$ alone, from week \nbSevenLaunchWeek{} onward, as a
proportional lift of \pc{\nbSevenLiftPct} on that market's untreated level:
%
\begin{equation}
  \Delta_t \;=\; 0.12 \,\big(c_0 + \lambda_{01} g_t + \lambda_{02} s_t + \lambda_{03} m_t\big)
  \cdot \mathbf 1[\,t \ge \nbSevenLaunchWeek\,],
  \qquad Y_{0t} \;\mathrel{+}=\; \Delta_t .
  \label{eq:app:geo_lift}
\end{equation}
%
Summed over the \nbSevenNPost{} post-launch weeks, the planted total lift is
\keur{\nbSevenTrueTotal}. \Cref{eq:app:geo_macro} is the single most consequential line in this
appendix: $m_t$ is a \emph{random walk}, so whatever share of it the donor weights fail to match
leaks into the estimator's error as a persistent offset rather than as fresh weekly noise --- and
that is why the Bayesian arm's iid likelihood underprices the uncertainty of the sum
(\cref{sec:sc:compare}). A world planted with iid errors would have hidden the book's sharpest
finding.

% =====================================================================================
\section{Staggered rollout}
\label{app:dgp:did}

\Cref{sec:did:dgm} draws \nbEightNStores{} stores, of which \nbEightNTreated{} receive the loyalty
app ($D_s = 1$), over \nbEightNWeeks{} weeks with the launch in week \nbEightLaunch:
%
\begin{equation}
  Y_{st} \;=\; \underbrace{\alpha_s}_{\text{store baseline}}
  \;+\; \underbrace{60 \sin\!\Big(\tfrac{2\pi t}{12}\Big)}_{\text{shared seasonal tide}}
  \;+\; \underbrace{400\, D_s\, \mathbf 1[\,t \ge 12\,]}_{\text{the app}}
  \;+\; \varepsilon_{st},
  \label{eq:app:did_y}
\end{equation}
%
with $\alpha_s \sim \mathcal N(1000,\ 150^{2})$ and
$\varepsilon_{st} \sim \mathcal N(0,\ 80^{2})$. The planted effect is \eur{\nbEightTrueEffect} per
store per week.

Two features are planted deliberately and both matter later. The seasonal term carries \emph{no
store-specific loading}: the tide lifts every store by the same euros in the same week, so
\textbf{parallel trends hold exactly, by construction} --- which is the only condition under which
a flat pre-trend can be recognised as a true negative rather than hoped to be one. And the store
baseline is \emph{large}: its standard deviation, \eur{\nbEightAlphaSd}, is nearly twice the weekly
noise, \eur{\nbEightNoiseSd}. That single number is what the difference-in-differences contrast
cancels, what an iid standard error charges the estimate for anyway, and what the pooled Bayesian
likelihood is forced to swallow into its residual variance --- which is why that chapter's
posterior comes out too \emph{wide}.

% =====================================================================================
\section{Threshold perks}
\label{app:dgp:rdd}

\Cref{sec:rd:dgm} draws \nbNineN{} customers with an annual spend $S_i$ --- the running variable
--- concentrated around the cutoff $c = \eur{\nbNineCutoff}$, and grants the Gold perk sharply:
%
\begin{align}
  S_i &\;\sim\; \mathcal N\!\left(100,\ 40^{2}\right)\ \text{clipped below at } 1,
        \qquad D_i \;=\; \mathbf 1\!\left[\,S_i \ge c\,\right], \label{eq:app:rd_s}\\
  p(S_i, D_i) &\;=\; \operatorname{clip}\!\Big(
      \underbrace{0.45}_{\text{baseline}}
      \;+\; \underbrace{0.0015\,(S_i - c)}_{\text{the confounder}}
      \;+\; \underbrace{\tau\, D_i}_{\text{the perk}},\ 0.02,\ 0.98\Big), \label{eq:app:rd_p}\\
  Y_i \mid S_i, D_i &\;\sim\; \text{Bernoulli}\big(p(S_i, D_i)\big),
      \qquad \tau \;=\; 0.14 . \label{eq:app:rd_y}
\end{align}
%
The slope $0.0015$ per euro is the \emph{smooth} part --- big spenders retain more anyway, and this
term knows nothing about any perk --- and the jump $\tau = \nbNineTrueJump$\,pp is the
\emph{causal} part, discontinuous by construction. Two properties that are facts here and
\emph{assumptions} on real data: the smooth part passes through the threshold without a kink, and
nobody manipulates $S$ to clear the tier.

% =====================================================================================
\section{The redesign}
\label{app:dgp:its}

\Cref{sec:its:dgm} produces \nbTenNDays{} days of a daily conversion rate, with the redesign
landing on day $t^{*} = \nbTenRedesignDay$ for everyone at once --- no control group exists, and
none ever will:
%
\begin{equation}
  Y_t \;=\; \underbrace{0.10 + 0.00015\,t}_{\text{level {+} drift}}
        \;+\; \underbrace{0.01 \sin\!\Big(\tfrac{2\pi t}{7}\Big)}_{\text{weekly cycle}}
        \;+\; \underbrace{0.03\;\mathbf 1[\,t \ge t^{*}\,]}_{\text{the redesign}}
        \;+\; \varepsilon_t,
  \qquad \varepsilon_t \sim \mathcal N\!\left(0,\ 0.006^{2}\right).
  \label{eq:app:its_y}
\end{equation}
%
The planted effect is a \emph{pure level step} of \nbTenTrueLift\,pp on a base of
\pc{\nbTenBaseRate}. The drift is the confounder, and it is planted small enough to be invisible on
a chart and large enough to dominate a before-and-after comparison: at $0.00015$ per day it
compounds across a pre-period into a bias that arithmetic alone can predict before any estimator
runs. The weekly cycle is the series' one genuinely autocorrelated feature, and it exists so that a
model which forgets it can be caught by a residual diagnostic.

% =====================================================================================
\section{Endogenous exposure}
\label{app:dgp:iv}

\Cref{sec:iv:dgm} draws \nbElevenN{} customers. Each carries an \textbf{unobserved} engagement
$U_i$ --- the intent the ad auction can see and the analyst cannot --- and receives a
\textbf{randomized} encouragement $Z_i$, a serving-priority lottery that makes the ad more likely
to reach them but is not itself an ad:
%
\begin{align}
  U_i &\sim \mathcal N(0, 1), \qquad Z_i \sim \text{Bernoulli}(\tfrac12), \label{eq:app:iv_u}\\
  X_i &\sim \text{Bernoulli}\Big(\sigma\big(0.3 + 1.1\,Z_i + 0.8\,U_i\big)\Big), \label{eq:app:iv_x}\\
  Y_i &= 50 \;+\; 15\,X_i \;+\; 12\,U_i \;+\; \varepsilon_i,
        \qquad \varepsilon_i \sim \mathcal N\!\left(0,\ 6^{2}\right). \label{eq:app:iv_y}
\end{align}
%
The planted effect of an exposure is \eur{\nbElevenTrue}. Everything the chapter is about is
visible in where $U$ appears: it enters \emph{both} \cref{eq:app:iv_x} and \cref{eq:app:iv_y}, and
that double appearance \emph{is} the self-selection. The instrument is the escape, and it is
defined by an \emph{absence} --- $Z$ shifts exposure in \cref{eq:app:iv_x} and appears nowhere in
\cref{eq:app:iv_y}. The exclusion restriction is not a property one can test; here, uniquely, it is
a property one can \emph{read}.

% =====================================================================================
\section{The worlds that are not simulated}
\label{app:dgp:real}

Four sections of this book have no data-generating model, because their data was generated by the
world. \Cref{sec:fnd:real} and \cref{sec:ctrl:real} use \citet{lalonde1986}'s job-training data, in
\citeauthor{dehejia1999}'s composite \citep{dehejia1999}. \Cref{sec:sc:real} uses the geo experiment
published with Google's geo-experiment methodology \citep{vaver2011, kerman2017}. \Cref{sec:iv:real}
uses Criteo's advertising experiment \citep{diemert2018}.

For these there is no planted truth, and the chapters do not pretend otherwise. What replaces it is
whatever the data itself can supply: a published experimental benchmark; an estimator that
randomization alone makes unbiased; a Wald estimate taken over the full file. That is a weaker
grade than a simulator gives --- and it is the strongest grade real data can ever give, which is
exactly why the simulated worlds above exist.
